% !TEX program = lualatex
% Transparencias de Fundamentos de las Matemáticas II
% Clase 04: Divisibilidad II

\documentclass[aspectratio=169,11pt]{beamer}

\usepackage{fontspec}
\usepackage[spanish,es-nodecimaldot]{babel}
\usepackage{amsmath,amssymb,mathtools}
\usepackage{booktabs}
\usepackage{csquotes}
\usepackage{xcolor}

\usetheme{Madrid}
\usecolortheme{default}
\setbeamertemplate{navigation symbols}{}
\setbeamertemplate{footline}[frame number]

\definecolor{FdMBlue}{RGB}{20,55,100}
\definecolor{FdMRed}{RGB}{130,30,30}
\definecolor{FdMGreen}{RGB}{25,100,60}

\setbeamercolor{structure}{fg=FdMBlue}
\setbeamercolor{title}{fg=white,bg=FdMBlue}
\setbeamercolor{frametitle}{fg=white,bg=FdMBlue}
\setbeamercolor{block title}{fg=white,bg=FdMBlue}
\setbeamercolor{block body}{bg=blue!3}
\setbeamercolor{alerted text}{fg=FdMRed}

\setmainfont{Latin Modern Roman}
\setsansfont{Latin Modern Sans}
\setmonofont{Latin Modern Mono}

\newcommand{\N}{\mathbb{N}}
\newcommand{\Z}{\mathbb{Z}}
\newcommand{\Q}{\mathbb{Q}}
\newcommand{\R}{\mathbb{R}}
\newcommand{\C}{\mathbb{C}}
\newcommand{\Pow}{\mathcal{P}}
\newcommand{\id}{\operatorname{id}}
\newcommand{\mcd}{\operatorname{mcd}}
\newcommand{\mcm}{\operatorname{mcm}}
\newcommand{\im}{\operatorname{Im}}

\title[FdM II -- Clase 04]{Divisibilidad II}
\subtitle{Máximo común divisor y mínimo común múltiplo}
\author{Antonio Falcó Montesinos}
\institute{Universidad CEU Cardenal Herrera}
\date{Fundamentos de las Matemáticas II}

\begin{document}

\begin{frame}
  \titlepage
\end{frame}

\begin{frame}{Objetivos de la clase}
\begin{itemize}
  \item Definir \(\mcd(a,b)\).
  \item Definir \(\mcm(a,b)\).
  \item Relacionar ambos conceptos.
\end{itemize}
\end{frame}

\begin{frame}{Mapa de la clase}
\tableofcontents
\end{frame}

\section{Máximo común divisor}

\begin{frame}{Máximo común divisor}
\begin{block}{Divisor común}
\begin{itemize}
  \item Un divisor común de \(a\) y \(b\) divide simultáneamente a ambos.
  \item El \(\mcd(a,b)\) es el mayor divisor común positivo.
  \item Se exige que \(a\) y \(b\) no sean ambos nulos.
\end{itemize}
\end{block}
\end{frame}

\begin{frame}{Máximo común divisor}
\begin{block}{Ejemplo}
\begin{itemize}
  \item Divisores positivos de \(18\): \(1,2,3,6,9,18\).
  \item Divisores positivos de \(30\): \(1,2,3,5,6,10,15,30\).
  \item Por tanto, \(\mcd(18,30)=6\).
\end{itemize}
\end{block}
\end{frame}

\begin{frame}{Máximo común divisor}
\begin{block}{Coprimos}
\begin{itemize}
  \item Dos enteros son coprimos si \(\mcd(a,b)=1\).
  \item Ejemplo: \(8\) y \(15\).
  \item La coprimalidad indica ausencia de factores comunes no triviales.
\end{itemize}
\end{block}
\end{frame}

\section{Mínimo común múltiplo}

\begin{frame}{Mínimo común múltiplo}
\begin{block}{Definición}
\begin{itemize}
  \item Un múltiplo común de \(a\) y \(b\) es divisible por ambos.
  \item \(\mcm(a,b)\) es el menor múltiplo común positivo.
  \item Se suele trabajar con \(a,b\neq 0\).
\end{itemize}
\end{block}
\end{frame}

\begin{frame}{Mínimo común múltiplo}
\begin{block}{Relación}
\begin{itemize}
  \item Para \(a,b\neq 0\), se cumple \(\mcd(a,b)\mcm(a,b)=|ab|\).
  \item Esta fórmula permite calcular uno conociendo el otro.
  \item Debe respetarse el valor absoluto.
\end{itemize}
\end{block}
\end{frame}

\section{Profundización}

\begin{frame}{Profundización}
\begin{block}{Interpretación}
\begin{itemize}
  \item El \(\mcd\) mide la parte común.
  \item El \(\mcm\) mide el primer punto común en los múltiplos positivos.
  \item Ambos resumen información de divisibilidad.
\end{itemize}
\end{block}
\end{frame}

\begin{frame}{Profundización}
\begin{block}{Ejemplo rápido}
\begin{itemize}
  \item \(\mcd(12,18)=6\).
  \item Como \(|12\cdot 18|=216\), entonces \(\mcm(12,18)=216/6=36\).
  \item Comprobación: \(36\) es múltiplo de \(12\) y de \(18\).
\end{itemize}
\end{block}
\end{frame}

\begin{frame}{Profundización}
\begin{block}{Error frecuente}
\begin{itemize}
  \item Usar el menor divisor común en lugar del mayor.
  \item Olvidar que el \(\mcm\) se define positivo.
  \item Aplicar la fórmula con \(a=0\) o \(b=0\) sin revisar hipótesis.
\end{itemize}
\end{block}
\end{frame}


\section{Detalle y práctica}

\begin{frame}{Detalle y práctica}
\begin{block}{División euclídea}
\begin{itemize}
  \item Si \(a\in\mathbb{Z}\) y \(b\in\mathbb{N}\), \(b>0\), existen únicos \(q,r\in\mathbb{Z}\) tales que \(a=bq+r\).
  \item El resto satisface \(0\leq r<b\).
  \item La unicidad del resto es la base de muchos argumentos posteriores.
\end{itemize}
\end{block}
\end{frame}

\begin{frame}{Detalle y práctica}
\begin{block}{Actividad breve}
\begin{itemize}
  \item Escribir la división euclídea de \(137\) entre \(11\).
  \item Escribir la división euclídea de \(-37\) entre \(5\) respetando \(0\leq r<5\).
  \item Comparar el resto euclídeo con el resto que da una calculadora elemental.
\end{itemize}
\end{block}
\end{frame}

\begin{frame}{Cierre}
\begin{itemize}
  \item El \(\mcd\) habla de divisores comunes.
  \item El \(\mcm\) habla de múltiplos comunes.
  \item La relación \(\mcd\cdot\mcm=|ab|\) conecta ambos.
\end{itemize}
\end{frame}

\begin{frame}{Trabajo recomendado}
\begin{itemize}
  \item Releer las secciones correspondientes del manual.
  \item Identificar definiciones, símbolos nuevos y enunciados principales.
  \item Preparar dudas y ejemplos para el seminario asociado.
\end{itemize}
\end{frame}

\end{document}
